\chapter{Visión del producto}

\section{Problema y oportunidad}

Las áreas administrativas reciben de forma continua boletas y comprobantes de pago de usuarios en formatos digitales heterogéneos con variaciones de diseño, legibilidad y contenido \cite{Zheng2023DocLevelIE} , \cite{Huang2021SROIE}. Esta diversidad obliga a transcribir y verificar manualmente campos clave, como el monto, fecha, número de operación, banco, entre otros. Lo que incrementa tiempos de conciliación, eleva el riesgo de errores y dificulta la trazabilidad ante auditorías\cite{Huang2021SROIE}. La literatura comenta que los documentos visualmente ricos combinan texto y disposición espacial, lo que plantea retos específicos y ha motivado la creación de conjuntos de datos y competencias orientadas a formularios y recibos.\cite{Cui2021}

En el campo de la investigación, múltiples trabajos han mostrado que modelar conjuntamente el contenido textual y la estructura de página mejora la extracción de información en documentos con layout variados. En especial, LayoutLM introdujo el preentrenamiento conjunto de texto y disposición con resultados competitivos en tareas de formularios y recibos, consolidando la importancia del layout en escenarios reales \cite{Xu2020LayoutLM}. En paralelo, se han publicado datasets para el análisis, DocLayNet, que busca generalizar más allá de dominios homogéneos y mitigar caídas de desempeño cuando los modelos se enfrentan a diferentes layouts que es algo frecuente que enviarían los usuarios \cite{Pfitzmann2022DocLayNet}. En conjunto, estos avances muestran que existe base científica y recursos públicos para abordar documentos con alta variabilidad en la disposición espacial de los datos.


Estos avances permiten enmarcar el caso de Cometa de manera directa, ya que esta startup ofrece a los colegios un portal de pagos para familias y, en este sentido, recibe carátulas y comprobantes con formatos diversos que deben validarse y estandarizarse antes de su uso operativo. Su plataforma opera en el contexto de la gestión de pagos escolares en México, facilitando que colegios particulares cobren colegiaturas y otros conceptos a través de una solución digital integrada con la pasarela de pagos Kushki \cite{cometa_onboarding}. Esto permite a los tutores, padres de familia, realizar pagos en línea de manera conveniente, revisando su historial desde el celular. Sin embargo, para que estos flujos funcionen correctamente, \textit{Cometa} debe validar y configurar las cuentas bancarias de los colegios donde se depositarán los fondos.


Actualmente, el proceso de validación de cuentas bancarias requiere que cada colegio proporcione una carátula bancaria, es decir, un documento oficial (por ejemplo, la portada de un estado de cuenta) donde constan los datos de la cuenta: nombre del titular, número de cuenta o CLABE interbancaria\footnote{La CLABE interbancaria (Clave Bancaria Estandarizada) es un número de 18 dígitos que identifica de forma única una cuenta bancaria en México y permite realizar transferencias electrónicas seguras entre bancos, según el Banco de México. Fuente: \url{https://www.banxico.org.mx/}}, banco, entre otros. El problema técnico radica en que dichas carátulas bancarias no poseen una estructura definida, ya que contienen logotipos, tablas, múltiples secciones y estilos de formato. Esta diversidad en presentación dificulta la extracción automática de la información relevante. Además, cada banco emite estados de cuenta con distintos diseños, y los usuarios pueden enviar documentos en PDF, imágenes escaneadas o fotos tomadas con el teléfono. 

Las soluciones tradicionales de OCR, , se quedan cortas cuando el contenido no sigue una plantilla uniforme. Por ello, hasta ahora la verificación suele hacerse manualmente: el equipo de soporte revisa la carátula bancaria y comprueba que los datos ingresados coincidan con el documento \cite{validacion_clabe}. Este proceso manual es lento y propenso a errores humanos, lo que crea un cuello de botella en la incorporación de nuevas cuentas.


En el contexto mexicano, el uso de inteligencia artificial para automatizar procesos documentales aún es incipiente. Algunas \textit{fintech} nacionales como \textit{Klar}, \textit{Konfío}, \textit{Kueski} o \textit{Clip} aplican algoritmos de IA en áreas como evaluación crediticia o detección de fraude, pero no existen herramientas orientadas a la \textit{extracción automática de información desde documentos bancarios institucionales} \cite{klar,konfio,kueski,clip}. Esto representa una oportunidad para que \textit{Cometa} desarrolle un sistema propio, ajustado a las necesidades locales y al flujo operativo del sector educativo \cite{cometa_inversion}.

A nivel internacional, servicios como \textit{Google Document AI}, \textit{Amazon Textract} o \textit{Microsoft Form Recognizer} han demostrado la efectividad del \textit{deep learning} en la comprensión de documentos \cite{layoutlm,doclaynet}. No obstante, su adopción en México enfrenta limitaciones: el costo por volumen de procesamiento, la dependencia de servidores externos y los requisitos de cumplimiento de la \textit{Ley Federal de Protección de Datos Personales en Posesión de los Particulares (LFPDPPP)} \cite{lfpdppp}. Además, estos modelos no han sido entrenados con los formatos de los bancos mexicanos, lo que reduce su precisión en el contexto local.

Ante este panorama, una solución desarrollada dentro de la infraestructura de \textit{Cometa} ofrece ventajas claras: control de la información, adecuación a las normativas nacionales y capacidad de personalizar los modelos para reconocer los diseños de bancos como BBVA, Banorte, Santander, HSBC México o Citibanamex \cite{cometa_inversion}.

El proyecto contempla la incorporación de modelos de lenguaje de gran escala (LLMs), como \textit{LLaMA 3} y \textit{Mistral}, para interpretar documentos no estructurados \cite{llama3,mistral}. Para garantizar la privacidad, estos modelos se entrenarían y ejecutarían en entornos locales o en nube privada dentro de México, evitando que los datos financieros salgan de la infraestructura de la empresa. Asimismo, se aplicarían técnicas de anonimización y preprocesamiento para reemplazar información sensible por etiquetas genéricas, manteniendo la confidencialidad sin afectar la comprensión del modelo.  

Asimismo, se propone aplicar técnicas de anonimización y preprocesamiento antes de enviar la información a la LLM: los nombres, números de cuenta y demás identificadores pueden sustituirse temporalmente por etiquetas genéricas. Esto permite aprovechar las capacidades del modelo sin comprometer la confidencialidad de los usuarios.  

En síntesis, esta extensión tecnológica representa una oportunidad para que \textit{Cometa} avance hacia un modelo operativo más ágil, seguro y escalable, alineado con su estrategia de crecimiento en el sector \textit{fintech} educativo mexicano \cite{cometa_inversion}.

En conjunto, la propuesta consolida una oportunidad estratégica para \textit{Cometa}: integrar inteligencia artificial avanzada en un proceso crítico de validación bancaria, bajo un marco de seguridad y cumplimiento con la normativa mexicana. Al desarrollar una solución propia basada en \textit{Document Understanding} y LLMs, la empresa no solo optimiza su operación, sino que fortalece su autonomía tecnológica y sienta las bases para futuros sistemas de automatización financiera en el sector \textit{fintech} educativo del país \cite{cometa_inversion}.

\section{Usuarios y casos de uso}

Los principales usuarios de este proceso son los administradores de los colegios que registran sus cuentas bancarias, y el equipo de soporte de Cometa, encargado de validar o completar registros cuando sea necesario. Para los colegios, contar con cuentas correctamente configuradas en Cometa es esencial: les permite recibir pagos sin contratiempos ni rechazos por datos erróneos. Gracias al sistema propuesto, al subir una carátula bancaria podrán ver los campos relevantes completarse automáticamente, reduciendo errores y acelerando la incorporación de nuevas cuentas.

Desde el lado de Cometa, el equipo de soporte también se beneficia al automatizar tareas que antes requerían validación manual. Esto permite escalar la operación a más colegios sin depender del crecimiento proporcional del equipo operativo, manteniendo la calidad del servicio con menos fricción \cite{cometa_inversion}.  

Actualmente, un colegio que se une a Cometa debe subir su carátula bancaria y esperar a que alguien del equipo revise y active la cuenta, lo que puede tomar horas o días. Con la solución propuesta, ese mismo administrador carga el documento y, en segundos, el sistema extrae los datos clave y los compara con los ingresados. Si hay coincidencia, la cuenta puede quedar validada casi en tiempo real. Si se detecta alguna discrepancia, el sistema puede alertar de inmediato. Esto se integra directamente en el flujo actual sin generar pasos adicionales.

Aunque los tutores (padres de familia) no intervienen en este flujo, su experiencia mejora indirectamente. Cada vez que realizan un pago, esperan que el dinero llegue correctamente a la cuenta del colegio. La automatización fortalece la confiabilidad general del sistema: los pagos realizados desde el front-end se enrutan correctamente en el back-end, y los colegios reciben los fondos sin fallas. Este tipo de eficiencia operativa responde directamente a los objetivos estratégicos de crecimiento de Cometa \cite{cometa_inversion}.

\section{Propuesta de valor y KPIs de éxito}

La propuesta de valor es automatizar la validación de cuentas bancarias mediante IA, logrando que la configuración sea más rápida, precisa y escalable. En una frase:  
\textit{“Cometa podrá extraer y verificar automáticamente la información de cuentas bancarias desde documentos cargados por los usuarios, eliminando errores manuales y acelerando cuando se podrán realizar los pagos escolares.”}

\subsection*{Objetivos SMART}
\begin{itemize}
  \item \textbf{Extracción automatizada precisa:} Desarrollar un prototipo de API capaz de identificar correctamente campos clave con al menos un 90\% de precisión para finales del Q2 2026 \cite{layoutlm}.
  \item \textbf{Reducción del tiempo de registro:} Reducir el tiempo promedio de validación a menos de 30 segundos.
  \item \textbf{Implementación piloto y ajuste:} Ejecutar un piloto con al menos cinco colegios antes de marzo de 2026, con un 80\% de carátulas procesadas sin intervención humana.
  \item \textbf{Cumplimiento regulatorio y de seguridad:} Asegurar conformidad con estándares bancarios y de protección de datos para Q4 2025.
\end{itemize}

\subsection*{Indicadores clave (KPIs)}
\begin{itemize}
  \item Precisión en extracción: $\geq$95\% en los campos principales.
  \item Tiempo promedio de registro: $<30$ segundos.
  \item Cuentas configuradas sin errores: $\geq$98\% sin reprocesamiento manual.
\end{itemize}

\noindent Estos indicadores permitirán monitorear el desempeño del sistema y su impacto operativo. Lograr estos objetivos fortalecerá la propuesta de valor de Cometa como plataforma fintech educativa confiable y eficiente.

\bibliographystyle{apalike}
